\documentclass[11pt]{ctexart}
\usepackage[a4paper,margin=1in]{geometry}
\usepackage{graphicx}
\usepackage{xcolor}
\usepackage{hyperref}
\definecolor{refgray}{gray}{0.45}
\title{\textbf{市场最新观点汇总}\\[0.4em]{\large 市场主线从需求侧叙事切换至供给侧再定价，AI资本开支与贸易保护重塑全球定价权，中国内需分化与能源冲击构成短期压力。}}
\author{KC桌面 自动生成}
\date{260527}
\begin{document}
\maketitle
\section*{一页摘要}
\begin{itemize}
  \item 全球AI资本开支持续超预期，2026年四大云服务商数据中心资本开支增速上调至+80\%，绝对金额增量超2500亿美元，市场仍低估其规模与持续性。
  \item 供给侧再定价成为多资产主线：中国煤炭供给因矿难后监管收紧面临结构性左移，西方钢铁利润率受贸易保护政策支撑，铁矿石边际成本系统性抬升。
  \item 中国内需呈现K型分化，四月数据走弱源于油价冲击与财政节奏放缓的双重挤压，但高端制造与出口韧性仍在。
  \item AI贸易规模已超能源贸易，半导体价格涨幅远超原油，嵌入AI供应链的经济体正经历贸易条件的结构性重估。
  \item 美国银行股估值折价反映资金结构问题而非基本面恶化，监管松绑正释放银行证券投资组合的结构性回补需求。
  \item 中国房地产市场二手房正成为新的定价锚，一线城市库存出清至临界点，但全国性复苏仍为时尚早。
\end{itemize}
\section{宏观与利率}
\textbf{全球利率环境正经历从'可预测锚'到'结构性波动'的范式转换，美联储决策哲学转变将系统性抬升利率波动率，而欧洲财政纪律回归正在挤出过去两年的宽松溢价。}

\begin{itemize}
  \item 新任美联储主席倾向'在房间里做决定'的决策模式，伯南克时代的渐进主义与前瞻指引将被取代，债券市场可预测性溢价将系统性下降，利率波动率将从事件性转为结构性。
  \item 债券市场对AI生产力故事仍持怀疑态度，2023年以来30多个主要AI模型发布日美国国债收益率几乎没有反应，市场认为AI是渐进式技术扩散而非颠覆性冲击。
  \item 欧洲主权债利差收窄至意大利与德国10年期国债利差约80个基点，但当前平静建立在脆弱支撑之上，欧洲委员会正在审查财政规则豁免条款，财政纪律回归将改变'宽松财政+央行兜底'的定价逻辑。
  \item 人民币中间价定价机制正经历从'政策锚'到'模型锚'的静默切换，模型预测值与实际中间价的持续偏离折射出政策层在汇率形成机制中逐步让渡空间给市场供求。
  \item 美国证券化信贷市场正经历供需两侧的同步修复，Basel终局提案修改释放了被压抑三年的银行证券投资组合容量，第一类和第二类银行CET1资本要求下降约4.8\%。
\end{itemize}
\begin{figure}[htbp]
\centering
\includegraphics[width=0.88\linewidth]{figures/fig_148.jpg}
\caption{Exhibit 1 - Exhibit 1: Our weekly composite index increased in the most recent week (+10\% w/w); the bottleneck scale remained at '2'; overall bottleneck levels remain well below peak congestion levels when scale was at '10' and now imply level}
\end{figure}
\begin{figure}[htbp]
\centering
\includegraphics[width=0.88\linewidth]{figures/fig_279.jpg}
\caption{Figure 1 - Figure 1: Trends in r-g across countries}
\end{figure}
\begin{figure}[htbp]
\centering
\includegraphics[width=0.88\linewidth]{figures/fig_283.jpg}
\caption{Figure 5 - Figure 5: EURO STOXX 600 and 30y BTP-bund spreads}
\end{figure}
\begin{figure}[htbp]
\centering
\includegraphics[width=0.88\linewidth]{figures/fig_288.jpg}
\caption{Figure 4 - Figure 4: Elevated inflation prints are coming but 1y1y inflation swaps at 2.60\% yoy in US and 2.16\% yoy in the euro area signal that markets are not panicking. Figure shows 1y inflation swap in 1y - Euro area and US (\%).}
\end{figure}
{\small\color{refgray}\textbf{References:} [R017] 债券市场正在为“后伯南克时代”的利率波动重新定价; [R015] 欧洲利差收窄并非警报解除，这轮放松的边界比市场想象的更窄; [R014] 人民币中间价模型正在发出一个被市场忽视的信号; [R011] 美国证券化信贷市场真正被低估的不是利差，而是银行需求的结构性回补}

\section{AI与科技硬件}
\textbf{AI基础设施正从'堆GPU'向'造工厂'范式转换，云数据中心资本开支的绝对金额跃迁正在重塑行业定价逻辑，而半导体供给侧紧张正从ABF载板向测试接口、功率半导体等更广泛环节蔓延。}

\begin{itemize}
  \item 2026年全球前四大美国云服务商数据中心资本开支增速从+63\%上调至+80\%，绝对金额增量超2500亿美元，2027年增速虽放缓至+50\%但绝对增量超2850亿美元，量级跃迁不能用'周期'简单定义。
  \item AI基础设施叙事正从单一GPU集群转向'AI工厂'级别的系统架构，NVIDIA正从服务器级设计走向机架级乃至Pod级设计，竞争门槛从'谁能造出最快的芯片'转向'谁能整合成可组合、可扩展的工厂级系统'。
  \item ABF载板供给紧张被市场系统性低估，AI芯片对载板的消耗量以代际为单位跳跃：Blackwell GPU消耗量为PC处理器的10.5倍，Rubin GPU达17倍，传统供需分析框架可能已失效。
  \item 功率半导体和测试接口的交期正在拉长，订单可见度提升，这轮供给紧张并非周期性补库，而是由结构性需求升级和产能约束共同驱动，相关公司议价能力和盈利持续性可能比市场预期更强。
  \item 中国半导体进口均价连续三个月同比超30\%，4月达1.0美元/颗，反映国内对高端AI、ADAS芯片的需求爆发，国产替代在高端领域渗透仍处早期阶段。
  \item SerDes速率升级正从渐进式改进转向结构性替代，224G是当前量产主力，但448G之后将触及物理极限，倒逼CPO商业化，AI网络正从'铜+光模块'方案转向CPO。
\end{itemize}
\begin{figure}[htbp]
\centering
\includegraphics[width=0.88\linewidth]{figures/fig_300.jpg}
\caption{Exhibit 4 - Exhibit 4: ABF substrate consumption by applications (by indexing PC processor to 1.0) There is a structural growth of ABF substrate demand across generations of AI CPU/GPU}
\end{figure}
\begin{figure}[htbp]
\centering
\includegraphics[width=0.88\linewidth]{figures/fig_321.jpg}
\caption{Figure 1 - Figure 1: Annual Data Center Capex from Top 4 U.S. CSPs}
\end{figure}
\begin{figure}[htbp]
\centering
\includegraphics[width=0.88\linewidth]{figures/fig_322.jpg}
\caption{Figure 2 - Figure 2: Annual Change in Data Center Capex from Top 4 U.S. CSPs (\$ in Billions Change)}
\end{figure}
\begin{figure}[htbp]
\centering
\includegraphics[width=0.88\linewidth]{figures/fig_301.jpg}
\caption{Exhibit 1 - Exhibit 1: Apr 2026 IC production YoY growth was \$+22.1\textbackslash{}\%\$ YoY, vs. \$+20.6\textbackslash{}\%\$ in Mar 2026 Semiconductor monthly data in Greater China Exhibit 2: IC production was +22.1\% YoY in Apr (vs. +20.6\% YoY in Mar) China's monthly IC produ}
\end{figure}
\begin{figure}[htbp]
\centering
\includegraphics[width=0.88\linewidth]{figures/fig_302.jpg}
\caption{Exhibit 3 - Exhibit 3: IC import value was +54.7\% YoY in Apr 2026 (vs. +53.7\% YoY in Mar 2026) China's monthly IC import value YoY}
\end{figure}
{\small\color{refgray}\textbf{References:} [R032] 市场仍在低估2026-2027年云数据中心资本开支的规模与持续性; [R021] AI基础设施正在经历一场从“堆GPU”到“造工厂”的范式转换; [R020] ABF载板供需紧张正在被市场低估，真正的拐点要到2028年以后; [R019] 功率半导体与测试接口的供给紧张，才是市场真正低估的结构性变量; [R023] 市场真正低估的不是需求，而是中国半导体进口均价的持续重构; [R042] 市场低估了SerDes对AI网络架构的再定价}

\section{能源与大宗商品}
\textbf{大宗商品定价逻辑正从需求侧切换至供给侧，中国煤炭供给因矿难后监管收紧面临结构性左移，铁矿石边际成本系统性抬升，西方钢铁利润率受贸易保护政策支撑，而AI贸易对能源的'对冲系数'被市场低估。}

\begin{itemize}
  \item 山西沁源煤矿事故后安全监管从'运动式'走向'制度化'，山西停产焦煤矿增至73座涉及产能7890万吨/年，中国煤炭供给的边际成本正在系统性上移，市场仍在用'需求疲软'逻辑定价。
  \item 铁矿石价格在供给增加、中国钢铁产量同比下降背景下逆势上涨，核心驱动力是边际成本的系统性抬升——中东冲突导致柴油、炸药等投入品价格上涨，海运费大幅攀升，长期实际价格从80美元/吨上调至90美元/吨。
  \item 全球钢铁市场正经历贸易保护重塑定价权，美国热轧卷板价格年初至今上涨近200美元/吨，欧洲德国热卷BOF价差涨至231美元/吨远高于周期均值，'需求决定价格'框架正在让位于'政策决定价格'新范式。
  \item 全球AI相关产品贸易规模已是能源贸易的近两倍，自2025年中以来存储芯片和半导体价格涨幅约为布伦特原油涨幅的八倍（455\%对55\%），嵌入AI供应链的经济体正经历贸易条件的结构性重估。
  \item 锂的供需拐点已经到来，BESS在锂需求中的占比将从2025年的约25\%跃升至2027/28年的约40\%，对应碳酸锂当量需求从约30万吨升至约90万吨，定价权之争才刚刚开始。
  \item 铜面临新的供给瓶颈，刚果金80-90\%的铜生产依赖从中东进口硫酸，若霍尔木兹海峡持续紧张，约0.8-1.0百万吨铜产量可能受影响。
\end{itemize}
\begin{figure}[htbp]
\centering
\includegraphics[width=0.88\linewidth]{figures/fig_007.jpg}
\caption{Exhibit 1 - Exhibit 1: ASP at coal mine pits Thermal coal prices at mine in Shanxi (Kcal5.5k)/Inner Mongolia (Kcal5.5k)/Shaanxi (Kcal5.8k) changed -0.4\%/+0.3\%/+0.9\% WoW to RMB692/574/652/t as of 22 May 2026}
\end{figure}
\begin{figure}[htbp]
\centering
\includegraphics[width=0.88\linewidth]{figures/fig_010.jpg}
\caption{Exhibit 4 - Exhibit 4: Liulin No. 4 hard coking coal price (incl. VAT)}
\end{figure}
\begin{figure}[htbp]
\centering
\includegraphics[width=0.88\linewidth]{figures/fig_092.jpg}
\caption{Figure 1 - Figure 1: HRC prices}
\end{figure}
\begin{figure}[htbp]
\centering
\includegraphics[width=0.88\linewidth]{figures/fig_094.jpg}
\caption{Figure 4 - Figure 4: German HRC vs BOF HRC cash costs (average producers)}
\end{figure}
\begin{figure}[htbp]
\centering
\includegraphics[width=0.88\linewidth]{figures/fig_285.jpg}
\caption{Figure 1 - Figure 1: Global AI enabling trade is larger than energy trade Trade as \% of global GDP by product group}
\end{figure}
\begin{figure}[htbp]
\centering
\includegraphics[width=0.88\linewidth]{figures/fig_198.jpg}
\caption{Figure 1 - Figure 1: China steel production - Large and Medium enterprises}
\end{figure}
\begin{figure}[htbp]
\centering
\includegraphics[width=0.88\linewidth]{figures/fig_199.jpg}
\caption{Figure 2 - Figure 2: Freight rates vs Oil price}
\end{figure}
{\small\color{refgray}\textbf{References:} [R002] 市场真正低估的不是需求，而是供给侧再定价的深度; [R008] 煤炭供给的再定价窗口已经打开，但市场仍在犹豫; [R009] 山西矿难真正冲击的不是煤炭，而是全球冶金煤的定价锚; [R010] 铁矿石的定价逻辑正在被成本曲线重新定义，而非需求; [R004] 全球钢铁市场的真正主线：贸易保护正在重塑定价权，而非需求复苏; [R016] 市场真正低估的不是能源冲击，而是AI贸易对能源的“对冲系数”; [R024] 山西停产的真实冲击不在产量，而在中国炼焦煤配煤体系的脆弱性; [R044] 中国焦煤供应冲击的真正风险，在煤矿之外; [R045] 市场真正低估的不是需求，而是供给侧的再定价}

\section{地产与消费}
\textbf{中国房地产市场二手房正成为新的定价锚，一线城市库存出清至临界点但全国性复苏仍为时尚早，经济呈现K型分化——内需偏弱但高端制造与出口韧性仍在。}

\begin{itemize}
  \item 新房与二手房成交分化加剧，截至5月24日当周50城新房周度成交同比下降9\%，而10城二手房周度成交同比增长45\%，年初至今新房累计同比-15\%而二手房累计同比+5\%，二手房正从'跟随'转向'主导'定价。
  \item 上海二手房挂牌量已从去年12万套降至8.5万套，隐含库存消化周期仅6个月，300万元以下低总价房源已率先出现价格上涨，卖家行为模式从'恐慌性抛售'切换至'理性博弈'。
  \item 四月经济数据全面走弱是油价冲击与财政节奏放缓的双重挤压，而非趋势性拐点，工业增加值环比下降1.0\%创三年最大单月跌幅，但高端制造业与出口部门韧性比市场预期更强。
  \item 消费动能持续疲软，乘用车零售和线上家电销售同比增速进一步走弱，政策效果递减，建筑链条依然低迷，水泥发运量和螺纹钢需求徘徊在多年低位。
  \item 全国性复苏是'K型分化'——核心城市改善型/高端产品表现更好，低线城市依然承压，政策端短期不会全国放松，政治局会议用'继续努力稳定'而非'止跌回稳'表述。
\end{itemize}
\begin{figure}[htbp]
\centering
\includegraphics[width=0.88\linewidth]{figures/fig_262.jpg}
\caption{Exhibit 1 - Exhibit 1: Primary weekly unit sales and four-week moving average – Total 50 cities}
\end{figure}
\begin{figure}[htbp]
\centering
\includegraphics[width=0.88\linewidth]{figures/fig_263.jpg}
\caption{Exhibit 1 - Exhibit 1: Container throughput growth held up in May MTD}
\end{figure}
\begin{figure}[htbp]
\centering
\includegraphics[width=0.88\linewidth]{figures/fig_275.jpg}
\caption{Exhibit 13 - Exhibit 13: Cement shipment is hovering at multi-year lows... Weekly Cement Shipments in China}
\end{figure}
\begin{figure}[htbp]
\centering
\includegraphics[width=0.88\linewidth]{figures/fig_276.jpg}
\caption{Exhibit 14 - Exhibit 14: ...as is rebar demand Weekly Rebar Demand}
\end{figure}
\begin{figure}[htbp]
\centering
\includegraphics[width=0.88\linewidth]{figures/fig_277.jpg}
\caption{Exhibit 15 - Exhibit 15: Primary housing sales remains lackluster Weekly Primary Sales in 50 Cities (units)}
\end{figure}
\begin{figure}[htbp]
\centering
\includegraphics[width=0.88\linewidth]{figures/fig_278.jpg}
\caption{Exhibit 16 - Exhibit 16: Secondary housing sales in major cities held up better Weekly secondary home sales in 10 cities (units)}
\end{figure}
{\small\color{refgray}\textbf{References:} [R012] 中国楼市正在发生一个被忽视的结构性变化：二手房正在成为新的定价锚; [R013] 市场真正低估的不是四月数据，而是两速经济中“出口韧性”的定价权; [R027] 市场真正低估的不是需求回暖，而是供给侧的结构性出清; [R033] 亚洲真正的分化：科技繁荣掩盖了中国需求的疲软}

\section{区域市场}
\textbf{亚洲经济活动呈现科技繁荣掩盖中国需求疲软的分化格局，日本半导体材料与设备受益于测试封装结构性升级和存储周期重启，韩国金融股防御性被低估。}

\begin{itemize}
  \item 亚洲二季度GDP面临非对称风险：中国4月内需数据全面低于预期，工业产出环比下降1.0\%创三年最大单月跌幅，而EMAX地区科技出口动能继续加速，新加坡科技出口年化增速已超200\%。
  \item 日本半导体材料股受益逻辑远不止于AI算力，半导体产业链价值增量正从前道制造向后道测试和先进封装迁移，日本在这一环节拥有全球最深的护城河，能否切入中国供应链已成为盈利分水岭。
  \item 日本半导体设备市场的真正拐点是存储周期再启动，DDR5现货价格从4.5美元跃升至34美元并维持高位，存储客户重新具备资本开支能力与意愿，设备需求'量价结构'正在改变。
  \item 韩国金融股防御性来源正发生结构性变化，从单纯股息率转向盈利韧性与股东回报政策同步强化，银行股估值修复进入'深水区'，保险股承保质量出现多年首次积极信号。
\end{itemize}
\begin{figure}[htbp]
\centering
\includegraphics[width=0.88\linewidth]{figures/fig_330.jpg}
\caption{Figure 1 - Figure 1: Asia\textbackslash{}* Nominal and Real Tech Exports}
\end{figure}
\begin{figure}[htbp]
\centering
\includegraphics[width=0.88\linewidth]{figures/fig_331.jpg}
\caption{Figure 2 - Figure 2: Asia\textbackslash{}* Tech Export Prices}
\end{figure}
\begin{figure}[htbp]
\centering
\includegraphics[width=0.88\linewidth]{figures/fig_384.jpg}
\caption{Figure 1 - Figure 1: China high-tech manufacturing IP}
\end{figure}
\begin{figure}[htbp]
\centering
\includegraphics[width=0.88\linewidth]{figures/fig_385.jpg}
\caption{Figure 2 - Figure 2: Housing policy and housing market activity}
\end{figure}
{\small\color{refgray}\textbf{References:} [R033] 亚洲真正的分化：科技繁荣掩盖了中国需求的疲软; [R039] 日本半导体材料股的真正机会不在AI算力，而在测试与封装的结构性升级; [R040] 半导体设备市场的真正拐点，不是AI需求，而是存储周期的再启动; [R025] 韩国金融股：市场波动中，防御性被低估了}

\section{金融与银行}
\textbf{美国银行股估值折价反映资金结构问题而非基本面恶化，监管松绑正释放银行证券投资组合的结构性回补需求；中国银行业净息差反弹可持续性被低估，产能过剩风险正在消退。}

\begin{itemize}
  \item 美国大型银行股跑输标普500指数670个基点，但基本面稳健——一季度贷款增速达6\%，投行业务收入同比增长30\%，交易收入增长17\%，盈利预期年内在上修，下跌本质是资金流向AI交易的结构性扭曲。
  \item 后金融危机时代长达15年的监管压制周期正在终结，Basel终局提案修改下第一类和第二类银行CET1资本要求下降约4.8\%，银行此前为应对2023年提案而额外持有的资本缓冲可被重新配置。
  \item 中国银行业净息差反弹来自贷款收益率企稳和资金成本持续下降的合力，并非信贷规模扩张驱动，而是在更健康的信贷环境下自然形成，反弹可持续性被市场低估。
  \item 中国高风险工业信贷占比已从2025年底的约10\%降至约7\%，工业产能扩张在2H25明显放缓，出口继续强劲，工业企业利润1Q26同比增19\%，产能过剩风险正在消退。
  \item 中国银行股估值折价反映的不是基本面问题，而是市场对结构性变化的定价尚未完成，理解这一点比预测下一次FOMC会议结果更重要。
\end{itemize}
\begin{figure}[htbp]
\centering
\includegraphics[width=0.88\linewidth]{figures/fig_045.jpg}
\caption{Exhibit 2 - Exhibit 2: Banks trade at a steep discount to the S\&P 500 Fwd P/E of S\&P 500 ex Mag7 vs BKX, KRX P/E}
\end{figure}
\begin{figure}[htbp]
\centering
\includegraphics[width=0.88\linewidth]{figures/fig_046.jpg}
\caption{Exhibit 3 - Exhibit 3: Bank P/E multiples, valuations between banks have not recovered to pre-GFC levels P/E valuation by year, average, and maximum P/E of BKX constituents that have existed 1990-present}
\end{figure}
\begin{figure}[htbp]
\centering
\includegraphics[width=0.88\linewidth]{figures/fig_052.jpg}
\caption{Exhibit 9 - Exhibit 9: Loan growth has rebounded on resilient GDP growth YoY GDP, loan growth 1987-present}
\end{figure}
\begin{figure}[htbp]
\centering
\includegraphics[width=0.88\linewidth]{figures/fig_054.jpg}
\caption{Exhibit 14 - Exhibit 14: Super regional banks expect 9\% RWA benefit, on average Super regional expected RWA benefit from capital proposals}
\end{figure}
\begin{figure}[htbp]
\centering
\includegraphics[width=0.88\linewidth]{figures/fig_402.jpg}
\caption{Exhibit 4 - Exhibit 4: Nominal GDP growth improved notably in 1Q26}
\end{figure}
\begin{figure}[htbp]
\centering
\includegraphics[width=0.88\linewidth]{figures/fig_405.jpg}
\caption{Exhibit 7 - Exhibit 7: PPI growth accelerated to 2.8\% YoY in April 2026}
\end{figure}
\begin{figure}[htbp]
\centering
\includegraphics[width=0.88\linewidth]{figures/fig_412.jpg}
\caption{Exhibit 19 - Exhibit 19: Manufacturing profit growth was 19.1\% YoY in 1Q26}
\end{figure}
{\small\color{refgray}\textbf{References:} [R003] 银行股真正的机会不在于降息，而在于一个被市场低估的结构性拐点; [R011] 美国证券化信贷市场真正被低估的不是利差，而是银行需求的结构性回补; [R038] 银行股真正的拐点不是净息差，而是产能过剩风险的消退}

\section{医疗健康与创新药}
\textbf{中国医疗健康产业正从'跟随式创新'走向'全球定价权'，AI药物发现已从概念验证进入可验证收入阶段，ASCO 2026数据表明中国biotech在多个适应症上开始挑战标准疗法。}

\begin{itemize}
  \item AI药物发现的生产率拐点已经到来，某AI制药公司平台平均仅需4.5个月即可生成一个新的临床前候选化合物，自2021年以来已累计产出约30个PCC，首款完全由AI设计的药物已完成IIa期临床试验。
  \item ASCO 2026摘要显示中国biotech在多个适应症上从'me-too'追赶阶段进入能够挑战甚至重新定义标准疗法的阶段，恒瑞医药在肝癌、结直肠癌等领域呈现最多关键数据。
  \item 美国医疗器械板块远期市盈率跌至14-15倍，相较标普500折价达30\%为三十年来最大，但下跌本质是资金被AI交易抽走的结构性扭曲，而非基本面恶化，行业有机增长在2024年达7\%左右。
  \item 跨国药企的'三角博弈'定价权正在重塑全球创新药竞争格局：美国SMID生物科技提供早期创新，日本大型药企提供商业化规模，中国凭借快速临床执行和成本效率成为不可忽视的资产来源地。
  \item 日本药企面临严峻专利悬崖和管线补充压力，武田口服TYK2抑制剂Zaso峰值销售指引定在30-60亿美元，口服银屑病药物市场预计从15\%份额扩大至30-40\%。
\end{itemize}
\begin{figure}[htbp]
\centering
\includegraphics[width=0.88\linewidth]{figures/fig_397.jpg}
\caption{Figure 1 - Figure 1: Large-Cap Organic Growth vs. NTM Earnings Multiples}
\end{figure}
\begin{figure}[htbp]
\centering
\includegraphics[width=0.88\linewidth]{figures/fig_398.jpg}
\caption{Figure 2 - Figure 2: MedTech P/E Premium (Discount) vs. the S\&P 500}
\end{figure}
{\small\color{refgray}\textbf{References:} [R026] ASCO 2026的真正信号：中国药企的临床数据正在从“跟随”转向“定义标准”; [R028] 中国医疗健康产业正在经历一次被低估的结构性转型：从“跟随式创新”走向“全球定价权”; [R029] AI制药的价值验证比市场预期的更近，但真正的胜负手在商业化落地; [R034] 市场真正低估的不是需求，而是供给侧的再定价; [R036] 市场真正低估的不是新药管线，而是跨国药企的“三角博弈”定价权}

\section{风险偏好与资金流向}
\textbf{市场风险偏好受AI叙事主导，但盈利增长已取代估值扩张成为市场上涨的唯一引擎，投机情绪指标远低于历史高点，市场尚未进入狂热阶段。}

\begin{itemize}
  \item 标普500指数2026年至今上涨约10\%，但同期前瞻12个月每股收益预期上调了15\%，市盈率反而下降4\%，过去两年指数全部涨幅均可归因于盈利增长而非估值扩张。
  \item 标普500指数有望在2026年底达到8000点，2026年每股收益预计达340美元同比增长24\%，2027年进一步增长13\%至385美元，盈利增长能否持续是市场真正风险所在。
  \item 投机情绪指标远低于历史高点，IPO规模相对温和，动量交易虽已到极端水平但投资者情绪指标仅0.3，市场还没到狂热阶段。
  \item 美国企业硬件OEM板块在联想财报后单日大涨13\%，但板块估值已来到历史性高位——与英伟达估值几乎持平，而历史上这个折价通常在60\%-70\%，市场可能对短期利好过度反应。
  \item 中国互联网板块出现短期情绪驱动的抛售与中长期结构性竞争力重塑的同步错位，美团近期跑输恒生科技指数约5.8个百分点主要源于散户资金流出而非基本面恶化。
\end{itemize}
\begin{figure}[htbp]
\centering
\includegraphics[width=0.88\linewidth]{figures/fig_164.jpg}
\caption{Exhibit 1 - Exhibit 1: Continued earnings growth should drive continued equity market upside}
\end{figure}
\begin{figure}[htbp]
\centering
\includegraphics[width=0.88\linewidth]{figures/fig_165.jpg}
\caption{Exhibit 2 - Exhibit 2: Earnings estimates have outpaced recent equity market appreciation}
\end{figure}
\begin{figure}[htbp]
\centering
\includegraphics[width=0.88\linewidth]{figures/fig_166.jpg}
\caption{Exhibit 3 - Exhibit 3: YTD change in S\&P 500 price, earnings estimates, and P/E multiple YTD change in S\&P 500 price, earnings, and valuation}
\end{figure}
\begin{figure}[htbp]
\centering
\includegraphics[width=0.88\linewidth]{figures/fig_169.jpg}
\caption{Exhibit 6 - Exhibit 6: Speculative Trading Indicator is well below recent and historical highs}
\end{figure}
\begin{figure}[htbp]
\centering
\includegraphics[width=0.88\linewidth]{figures/fig_170.jpg}
\caption{Exhibit 7 - Exhibit 7: IPO activity has been relatively modest Number of US IPOs}
\end{figure}
\begin{figure}[htbp]
\centering
\includegraphics[width=0.88\linewidth]{figures/fig_432.jpg}
\caption{Exhibit 1 - Exhibit 1: Enterprise Hardware OEMs trade at a median 93\% relative valuation to NVDA today, well above historical levels.}
\end{figure}
\begin{figure}[htbp]
\centering
\includegraphics[width=0.88\linewidth]{figures/fig_433.jpg}
\caption{Exhibit 2 - Exhibit 2: US OEMs now trade at their highest relative valuation to the Taiwanese ODMs since 2018.}
\end{figure}
{\small\color{refgray}\textbf{References:} [R007] 市场真正低估的不是AI泡沫，而是盈利增长正在重塑估值叙事; [R041] 市场真正低估的不是需求，而是供给侧的再定价; [R043] 中国互联网的战术窗口与结构性隐忧}

\section*{结语}
综合来看，当前市场主线正从需求侧叙事切换至供给侧再定价，AI资本开支的绝对金额跃迁与贸易保护政策正在重塑全球定价权，而中国内需分化与能源冲击构成短期压力。投资者需关注供给侧结构性收紧对资产定价的深远影响，以及利率波动率从事件性向结构性转变带来的固定收益组合管理挑战。
\vfill
{\small\color{refgray}Personal reading notes and learning share only. Not investment advice.}
\end{document}
